\section{Actual analytic superlevels of coherent ideal-order tuples}

The analytic semicontinuity used in Bierstone--Milman\cite{bierstone-milman97-invent-math},
Section 13, is stronger than topological upper semicontinuity. Here the
superlevels are constructed from genuine coherent ideal orders, including
infinite order, before any dimensional component is selected. This proves
the finite global ideal-order tuple layer. It does not identify an arbitrary
enhanced resolution invariant with such a tuple.

\begin{definition}[Constructed lexicographic order locus]
\label{hir:actual-lex-order-locus}
\lean{Hironaka.HironakaOrderTupleAnalyticSuperlevels.orderTupleSuperlevel}
\leanok
Fix a closed smooth submanifold $N$ and a finite list of coherent ideals
with natural thresholds $L=((I_1,\mu_1),\ldots,(I_s,\mu_s))$.
Define $A_\varnothing=N$ and recursively
\[
 A_{(I,\mu)::L}=\operatorname{cosupp}_N(I,\mu+1)
 \cup\bigl(\operatorname{cosupp}_N(I,\mu)\cap A_L\bigr).
\]
\end{definition}
\begin{proof}
This finite recursion uses actual order cosupports, with the zero threshold
giving all of $N$. No superlevel analyticity or extension certificate is
supplied as data.
\end{proof}

\begin{lemma}[The actual lexicographic dictionary]
\label{hir:actual-lex-order-dictionary}
\uses{hir:actual-lex-order-locus}
\lean{Hironaka.HironakaOrderTupleAnalyticSuperlevels.mem_orderTupleSuperlevel_iff}
\leanok
The constructed locus equals
\[
 \{x\in N:(\mu_1,\ldots,\mu_s)\le_{\rm lex}
 (\operatorname{ord}_x(I_1|_N),\ldots,\operatorname{ord}_x(I_s|_N))\},
\]
where orders lie in $\mathbb N\cup\{\infty\}$ and the thresholds are finite.
\end{lemma}
\begin{proof}
For an empty list both sides are $N$. For a nonempty list, the actual
order theorem identifies cosupport membership with the inequality on
the first extended-natural order. Strictly exceeding a finite threshold
$\mu$ is equivalent to order at least $\mu+1$, also for infinite order.
Lexicographic comparison of two nonempty lists says either the first
entry strictly exceeds its threshold, or it equals its threshold and
the tail comparison holds. Equivalently it says that the first entry
strictly exceeds its threshold, or is at least its threshold and the
tail comparison holds. The recursive formula and induction give this
exact equivalence, with membership in $N$ in either branch.
\end{proof}

\begin{theorem}[Produced global analytic tuple superlevels]
\label{hir:actual-lex-order-analytic-superlevels}
\uses{hir:actual-lex-order-dictionary}
\lean{Hironaka.HironakaOrderTupleAnalyticSuperlevels.isComplexAnalyticSubset_orderTupleSuperlevel}
\leanok
Every constructed order-tuple superlevel is a global analytic subset,
without compactness or nonvanishing assumptions on the ideals.
\end{theorem}
\begin{proof}
The closed smooth $N$ is analytic. At a positive threshold $r$ the actual
cosupport is $N\cap V(\mathcal D_N^{r-1} I)$, with a coherent derivative
ideal, so it is analytic. At threshold zero, actual holomorphicity makes
cosupport exactly $N$, so it is analytic as well. Global analytic sets
are closed under finite intersection and union. Induction through the
displayed recursion proves the assertion. A zero ideal, hence infinite
order, causes no failure of derivative coherence or of this argument.
\end{proof}

\begin{theorem}[Actual global pure extension of a regular tuple trace]
\label{hir:actual-lex-order-global-pure-extension}
\uses{hir:actual-lex-order-analytic-superlevels,hir:dimension-part-regular-trace,hir:dimension-part-global-component-union}
\lean{Hironaka.HironakaOrderTupleAnalyticSuperlevels.exists_pure_global_extension_of_regular_tuple_superlevel}
\leanok
If the actual order-tuple superlevel $A$ is regular of dimension $q$ at
every point of its trace on an open $U$, there is an internally produced
global analytic pure-$q$ set $T\subset A$, with $T\cap U=A\cap U$,
regular of dimension $q$ on that trace. It equals the union of $A$'s
actual global components of geometric pure dimension $q$.
\end{theorem}
\begin{proof}
Construct $T=A^{(q)}=\overline{\operatorname{Reg}_q A}$ using the preceding
actual global analytic set. Dimension-part analyticity, geometric purity,
closedness and the actual global-component identification give all global
assertions. The regular-trace lemma gives exact trace equality, and a
regular coordinate neighbourhood in $A$ is also a regular neighbourhood
in its dimension part. Singularities of $T$ outside $U$ are not excluded:
they remain the separate auxiliary-chain smoothification problem of
Bierstone--Milman, Lemma 13.7.
\end{proof}

\begin{theorem}[The actual canonical invariant's first-entry superlevels]
\label{hir:canonical-first-entry-analytic-superlevel}
\uses{hir:actual-lex-order-analytic-superlevels,hir:actual-lex-order-dictionary,hir:canonical-centres}
\lean{Hironaka.HironakaOrderTupleAnalyticSuperlevels.canonical_first_superlevel_eq_orderTupleSuperlevel,Hironaka.HironakaOrderTupleAnalyticSuperlevels.isComplexAnalyticSubset_canonical_first_superlevel}
\leanok
For positive ambient dimension, let $\mathfrak m$ be an actual marked
ideal in the invariant domain and let $T$ be any genuine admissible test
tower. The locus in $N_T$ on which the first entry of the constructed
canonical invariant in positive induction dimension is at least $\mu$
is a global analytic subset. It equals the internally constructed
one-entry order-tuple locus for the actual nonmonomial ideal
$\mathcal N_{N_T}(I_T)$ and threshold $\mu$.
\end{theorem}
\begin{proof}
The actual transform remains in the invariant domain: its ideal is
coherent, its carrier is closed smooth, the genuine transformed boundary
is SNC and transverse, and the ideal is nonzero along its carrier.
The nonmonomial coherence theorem, proved by actual local monomial
division, therefore gives coherence of $\mathcal N_{N_T}(I_T)$. At every
point of $N_T$, its finite natural order used by the constructed invariant
equals its extended-natural ideal order. Indeed membership in every
actual order cosupport characterizes both orders; equality follows from
all finite threshold comparisons.

By the definitions of the genuine dimension-inductive invariant, its
first entry is exactly this nonmonomial order. Apply the actual one-entry
lexicographic dictionary to identify the stated locus with the produced
tuple locus, then apply its global analyticity theorem. This selects no
tuple as additional caller data and uses no analyticity assumption on the
canonical invariant. The higher companion/history entries, the monomial
second-entry superlevels, and the Hilbert--Samuel auxiliary-chain invariant
of Lemma 13.7 remain separate obligations.
\end{proof}
